% ==========================================
% 6. Discussion
% ==========================================
%
% Status: common, written together after the results are complete. This is
% the section where the five experiments are combined into one coherent
% analysis; it is not a juxtaposition of five separate discussions.
%
% Principle: we interpret the results, answer the research questions and
% acknowledge the limits. We introduce no new measurement here.
% We distinguish what the data show (evidence) from what we suppose
% (hypothesis for explanation), and we word each accordingly.

\section{Discussion}
\label{sec:discussion}

% ------------------------------------------
% Short framing paragraph
% ------------------------------------------
% TODO [LAST]: We announce the structure: answers to the research questions,
%              cross-experiment analysis, comparison with the literature,
%              limitations, implications.

\subsection{Answers to the Research Questions}
\label{sec:discussion-rq}
% TODO [LAST]: For each research question (RQ1, RQ2, ...), we state whether the
%              results support, partially support or contradict the hypothesis,
%              with the evidence (reference to results, figures, tables).
% TODO [LAST]: We keep the original hypotheses as written in
%              Section~\ref{sec:rq-hypotheses}; where the results differ, we
%              discuss the gap and do not rewrite the hypothesis.

\subsection{Cross-Experiment Analysis}
\label{sec:discussion-cross}
% TODO [LAST]: We compare the experiments with each other: common trends,
%              contradictions, complementary findings, dependencies between factors.
% TODO [LAST]: We explain differences between experiments by their protocol or
%              environment (Section~\ref{sec:methodology}) when they are
%              relevant.

\subsection{Unexpected Results}
\label{sec:discussion-unexpected}
% TODO [LAST]: We discuss the results that contradicted our expectations or
%              that remain unexplained, and propose explanations clearly
%              labelled as hypotheses, with a way to test them.

\subsection{Comparison with the State of the Art}
\label{sec:discussion-sota}
% TODO [LAST]: We compare our findings with the works of
%              Section~\ref{sec:state-of-the-art}, and state whether they confirm,
%              extend or contradict them. We check that the comparison is fair
%              (same metrics, same conditions).

\subsection{Limitations and Threats to Validity}
\label{sec:discussion-limits}
% TODO [LAST]: Internal validity: we list factors that may have biased the
%              measurements (environment, implementation, measurement error).
% TODO [LAST]: External validity: we state how far the results generalize (data,
%              hardware, scale).
% TODO [LAST]: Construct validity: we discuss whether the chosen metrics really
%              measure what the research questions ask.
% TODO [LAST]: We state the limits of the study without minimizing them, and
%              link each one to the conclusions it affects.

\subsection{Implications}
\label{sec:discussion-implications}
% TODO [LAST]: We describe the practical and scientific consequences of the
%              findings (for Splinter3D and for the field), in proportion to the
%              evidence.

% TODO [CHECK]: No result is presented here that does not appear in
%               Section~\ref{sec:results}.
% TODO [CHECK]: Every research question stated in Section~\ref{sec:rq-hypotheses}
%               receives an answer in Section~\ref{sec:discussion-rq}.
